\section{Certified diagonalization and the incompleteness chains}
\label{sec:p05-inc}

The following reconstruction uses the governing arithmetic language,
constructor numbering, and classical \(\mathsf{ND}_{\mathsf{FOL}}\) calculus of
\cref{sec:arithmetic-coding-04}. A computably enumerable theory is a
computably enumerable set of sentence axioms, not an axiom-membership
decider. All internal finite comparisons use the declared abbreviation
\(a\leq b\defeq\exists z\,(z+a=b)\), with its witness on the left of
addition. No induction schema is added to \(\mathsf Q\).

The proofs isolate their arithmetic input precisely. They reconstruct the
modern \(\mathsf Q\)-based versions; they do not identify that base theory
or numbering with G\"odel's original system \(P\)
\cite[Satz V--VI, pp.~186--190]{Godel1931} or with an unspecified historical
proof predicate \cite{Rosser1936}.

\begin{hbdefinition}{def:p05-inc-interface}{The pointwise certificate interface}
Fix \(x_0\) as the distinguished substitution variable. Let \(d(n)\)
be the code obtained by substituting \(\bar n\) for \(x_0\) in a formula
with code \(n\) and no other free variables, and a fixed sentence code
otherwise. Let \(\operatorname{neg}(n)\) be the code of the negation of
the formula with code \(n\), with a fixed default on malformed input.
These are numerical functions, not new symbols of \(\Lang_A\).

The diagonal certificate interface consists of a bounded formula
\(\delta_d(x,y,\vec s)\), the literal prefix formula
\(D(x,y)\defeq\exists\vec s\,\delta_d(x,y,\vec s)\), and a formula
\(N(x,y)\), such that, for each standard \(n\),
\[
 \mathsf Q\derives\forall y\,
       (D(\bar n,y)\leftrightarrow y=\overline{d(n)}),
 \qquad
 \mathsf Q\derives\forall y\,
       (N(\bar n,y)\leftrightarrow y=\overline{\operatorname{neg}(n)}).
\]
An empty certificate tuple is allowed. These are separately supplied,
numeral-by-numeral derivations, not the assertion that \(\mathsf Q\)
proves uniform totality, global graph determinism, or induction.

For each fixed c.e. \(T\supseteq\mathsf Q\), the proof interface consists
of its actual certified numerical relation \(\mathsf{Proof}_T(p,a)\)
from \cref{thm:arithmetization-syntax}, and a chosen \(\Sigma^0_1\)
formula \(P_T(p,a)\) satisfying both numeralwise clauses:
\[
 \mathsf{Proof}_T(p,a)\Rightarrow
       \mathsf Q\derives P_T(\bar p,\bar a),\qquad
 \neg\mathsf{Proof}_T(p,a)\Rightarrow
       \mathsf Q\derives\neg P_T(\bar p,\bar a).
\]
It must also define that relation in \(\mathcal N\). Put
\(\operatorname{Pr}_T(a)\defeq\exists p\,P_T(p,a)\).
When this is the predicate selected in \cref{def:fixed-proof-predicate},
\(P_T\) and \(\operatorname{Pr}_T\) are that definition's
\(\mathsf{Proof}_T\) and \(\mathsf{Prov}_T\). An alternative selected
formula is not thereby proved equivalent to it inside \(\mathsf Q\).
\end{hbdefinition}

These arithmetic inputs are supplied, rather than assumed as a new
representation principle. Select
\(D=\gamma_{d_{\rm diag}}\) and \(N=\gamma_{d_{\rm neg}}\)
from \cref{thm:p05-rep-diagonal-graph}; their numerical functions and
defaults are exactly those above. The bounded matrix and unrestricted
output graph proofs are given by
\cref{thm:p05-rep-trace,thm:p05-rep-strong-graph}. Select
\(P_T=\mathsf{Proof}^{*}_T\) and
\(\operatorname{Pr}_T=\mathsf{Prov}^{*}_T\), with the fixed
c.e. axiom-certificate convention of
\cref{def:p05-rep-certified-presentation}. Then
\cref{thm:p05-rep-certified-proof} supplies both numeralwise clauses,
standard-model accuracy and the literal bounded-matrix prefix. Thus the
following explicit interface hypotheses have an established instance for
every c.e. \(T\supseteq\mathsf Q\). Keeping them displayed also states
exactly which facts a different predicate selection would have to supply.
For this selection the certificate tuple for \(D\) has one coordinate.

\begin{hblemma}{thm:p05-inc-diagonal-coding}{proved=annex:p5.inc.diagonal-coding}{The actual substitution and negation codes}
The functions \(d\) and \(\operatorname{neg}\) just defined are primitive
recursive in the governing numbering. If \(\beta\) has at most free
variable \(x_0\), \(e=\ulcorner\beta\urcorner\), and
\(G=\beta[\bar e/x_0]\), then \(G\) is a sentence and
\(d(e)=\ulcorner G\urcorner\). The code is that of the fully expanded
formula, including the expansions of marked bounded quantifiers.
\end{hblemma}

\begin{hblemma}{thm:p05-inc-diagonal}{proved=annex:p5.inc.diagonal}{Pointwise graphs give the ordinary fixed point}
Assume the diagonal certificate interface of
\cref{def:p05-inc-interface}. For every formula \(\theta(x_0)\) with
at most that free variable one can construct a sentence \(G\) with
\[
 \mathsf Q\derives
 G\leftrightarrow\theta(\overline{\ulcorner G\urcorner}).
\]
Only the graph derivation at the single standard code used by the
construction is needed. This result makes no complexity claim about its
ordinary existential construction.
\end{hblemma}

\begin{hbtheorem}{thm:p05-inc-pi1-fixedpoint}{proved=annex:p5.inc.pi1-fixedpoint}{Complexity preservation without a collection axiom}
Assume the same diagonal certificate interface. If \(\theta(x_0)\)
is \(\Pi^0_1\) in the active syntactic hierarchy, its fixed point can
be chosen \(\Pi^0_1\). Explicitly, after fresh renaming and substitution,
put
\[
 \beta(x_0)\defeq\forall y\,\forall\vec s\,
       (\delta_d(x_0,y,\vec s)\to\theta(y)),\qquad
 e=\ulcorner\beta\urcorner,\qquad G=\beta(\bar e).
\]
Then \(d(e)=\ulcorner G\urcorner\) and the biconditional is proved
in \(\mathsf Q\). If \(\theta(v)\) is literally
\(\forall\vec u\,\eta(v,\vec u)\), with bounded \(\eta\), one
may instead use the literal prefix construction
\[
 \beta(x_0)\defeq\forall y\,\forall\vec s\,\forall\vec u\,
       (\delta_d(x_0,y,\vec s)\to\eta(y,\vec u)).
\]
All quantified variables are distinct and absent from the other formula's
free variables and bounds. For the strict prefix-form convention, take
the displayed bounded matrix's fixed negation-normal-form conversion
before computing \(e\). The resulting formula is already prefix form;
neither version invokes a \(\mathsf Q\)-provable bounded-collection
principle or merely replaces a formula by an \(\mathcal N\)-equivalent
prenex formula.
\end{hbtheorem}

\begin{hblemma}{thm:p05-inc-proof-facts}{proved=annex:p5.inc.proof-facts}{Proof evidence and the stronger consistency hypotheses}
Assume the proof interface for \(T\). For every sentence \(\sigma\),
\[
 T\derives\sigma\Longrightarrow
 T\derives\operatorname{Pr}_T(\overline{\ulcorner\sigma\urcorner}),
 \qquad
 \mathcal N\entails\operatorname{Pr}_T(\overline{\ulcorner\sigma\urcorner})
 \Longleftrightarrow T\derives\sigma.
\]
For a \(T\) extending \(\mathsf Q\), both \(\Sigma^0_1\)-soundness
and omega-consistency imply consistency. These facts require no internal
derivability condition for modus ponens or positive introspection.
\end{hblemma}

\begin{hbtheorem}{thm:p05-inc-godel-chain}{proved=annex:p5.inc.godel-chain}{The three precise G\"odel branches}
Let \(T\supseteq\mathsf Q\) be c.e. and assume the certificate
interfaces above. Construct \(G_T\) by
\cref{thm:p05-inc-pi1-fixedpoint} for
\(\theta(v)=\neg\operatorname{Pr}_T(v)\), and put
\(g=\ulcorner G_T\urcorner\). Thus \(G_T\) is \(\Pi^0_1\) and
\(\mathsf Q\derives G_T\leftrightarrow\neg\operatorname{Pr}_T(\bar g)\).
Then:
\begin{enumerate}
\item consistency of \(T\) implies \(T\notderives G_T\), and
      \(\mathcal N\entails G_T\);
\item \(\Sigma^0_1\)-soundness of \(T\) implies
      \(T\notderives\neg G_T\);
\item omega-consistency of \(T\) implies
      \(T\notderives\neg G_T\).
\end{enumerate}
If the selected proof formula is literally
\(P_T(p,a)=\exists\vec w\,\kappa_T(p,a,\vec w)\), bounded
\(\kappa_T\), the prefix construction for \(G_T\) is explicitly
\[
 \forall y\,\forall\vec s\,\forall p\,\forall\vec w\,
 (\delta_d(\bar e,y,\vec s)\to
                         \neg\kappa_T(p,y,\vec w)).
\]
Here \(e\) is the code of the unary prefix \(\beta\) built for
this selected \(\kappa_T\), with its bounded matrix converted to
NNF before coding, as specified in
\cref{thm:p05-inc-pi1-fixedpoint}. It is not the axiom-enumerator
index \(e_T\).
This is an additional formula-selection clause, not a proof that every
\(\Sigma^0_1\) representative is \(\mathsf Q\)-equivalent to prefix
form. Under either hypothesis (2) or (3), \(G_T\) witnesses deductive
incompleteness. Consistency alone is not used to infer its second branch.
\end{hbtheorem}

\begin{hbtheorem}{thm:p05-inc-rosser-chain}{proved=annex:p5.inc.rosser-chain}{Rosser's two finite comparison branches}
Let \(T\supseteq\mathsf Q\) be c.e., consistent, and equipped with
the same interfaces. One can construct a sentence \(R_T\), with
\(r=\ulcorner R_T\urcorner\) and
\(r^-=\ulcorner\neg R_T\urcorner=\operatorname{neg}(r)\), such that
\[
 \mathsf Q\derives R_T\leftrightarrow
 \forall y\bigl(P_T(y,\bar r)\to
                   \exists z\leq y\,P_T(z,\overline{r^-})\bigr).
\]
Neither \(R_T\) nor \(\neg R_T\) is derivable in \(T\).
The finite comparison uses \cref{lem:q-bounded-cases,lem:q-dichotomy};
it does not require \(T\) to be omega-consistent or sound. With a
\(\Sigma^0_1\) proof formula, the displayed Rosser matrix is not
asserted bounded and this construction is not advertised as a
\(\Pi^0_1\) Rosser sentence.
\end{hbtheorem}

\begin{hblemma}{thm:p05-inc-finite-extension}{proved=annex:p5.inc.finite-extension}{Finite extension consistency is a theoremhood query}
For every finite list \(D\) of sentences and every theory
\(T\supseteq\mathsf Q\), let \(C_D\) be their conjunction, or
\(0=0\) for an empty list. Then
\[
 T\cup D\text{ is inconsistent}
 \quad\Longleftrightarrow\quad T\derives\neg C_D.
\]
If \(T\cup D\) is consistent, at least one of its extensions by
\(\sigma\) and by \(\neg\sigma\) is consistent. The statement is
about the actual classical natural-deduction calculus, including its
nonvacuous discharge and quantifier side conditions.
\end{hblemma}

\begin{hbtheorem}{thm:p05-inc-theoremhood-chain}{proved=annex:p5.inc.theoremhood-chain}{Undecidability of theoremhood, not merely an undecided sentence}
For every consistent c.e. axiom theory \(T\supseteq\mathsf Q\),
equipped with the interfaces above for c.e. extensions, the set
\[
 \{\ulcorner\sigma\urcorner:\sigma\text{ is a sentence and }
                                    T\derives\sigma\}
\]
is computably enumerable and undecidable. If its membership were decidable,
there would be an effective sequence of finite sentence additions whose
union with \(T\) is a c.e., consistent, deductively complete extension.
The deductive closure of that extension would itself be decidable by
parallel proof search. Rosser's theorem forbids the extension already
at its c.e., consistent, complete stage.
\end{hbtheorem}

